\section{What this establishes and what it does not}
\paragraph{A function-prior statement.}
Theorem~\ref{thm:prior} bounds actual prior probabilities of function neighborhoods. Theorem~\ref{thm:energy-characterization} identifies the infinite-width minimum quadratic energy exactly with twice minimal path variation. These are stronger claims than calling a sample-complexity upper bound a prior. They still do not assert that every pair of functions is ranked exactly by path variation or that posterior computation is efficient.

\paragraph{The support condition changes parameter preferences.}
Absolute spectral nonexpansiveness is stronger than ordinary operator-norm nonexpansiveness: signed cancellations cannot hide amplification. Yet Lemma~\ref{lem:flow} shows that at a fixed width and depth it removes no ReLU functions. What changes is the parameterization and its Gaussian weight. General small networks may require large path variation, so this fact alone gives no unrestricted polynomial dependence jointly on teacher width and depth.

\paragraph{The Gaussian scaling is explicit.}
The first layer and readout have common precision $N(d+1)+(H-1)N^{\gamma+1}$, while internal coordinates have variance $N^{-\gamma}$. The likelihood is calibrated to the endpoint precision. This is a particular family of layerwise priors; the theorem does not transfer automatically to the more usual tuple with identical normalized strength at all layers. Its internal exponent can be any $\gamma\ge1$. More generally, internal precisions $\alpha_\ell>0$ may vary by layer: choosing $\omega=N(d+1)+N\sum_{\ell=2}^H\max\{N,\alpha_\ell\}$ and internal variances $1/\alpha_\ell$ gives the same results. Total dimension divided by $\omega$ is at most one, and the cloned internal mean cost is at most $C/N$; all other proofs are unchanged. This extension concerns calibrated endpoint priors, not an exhaustive claim about arbitrary unchanged scaling tuples.

\paragraph{Dimension and learner width.}
The statistical rate is dimension-free on the Euclidean unit ball, at every learner width above the comparison width. The architecture-dependent precision includes the first-layer parameter volume, so no additional lower bound on width in terms of input dimension is needed. The prior is specified separately at each width. If one starts with a cube and coordinatewise coefficient caps, normalizing the input and its constant coordinate changes first-row Euclidean norms; those factors must be counted in the teacher path variation. The theorem does not erase that distinction between norm and coordinate bounds.

\paragraph{Noise and approximation.}
Labels need not be generated by a network or by Gaussian noise. The likelihood is selected from the observed free energy and automatically gives a faster bound when the comparison risk is small. Prediction averages raw outputs. The result supplies neither an efficient free-energy calculation nor a gradient-descent guarantee.

\paragraph{Verification status.}
The arguments are conventional proofs using the published smooth-loss theorem with its assumptions stated explicitly. They have not been formalized in Lean.
